\SourcePage{376}{381}
\section{Breit equation}

As is well known, in classical electrodynamics a system of interacting particles can be described with the help of a Lagrangian function dependent only on the coordinates and velocities of the particles themselves and correct to terms $\sim 1/c^2$ (see~II, \S\,65). This circumstance is connected with the fact that radiation appears only as an effect of order~$1/c^2$.

In quantum theory this situation corresponds to the possibility of description of the system by a Schr\"odinger equation, with perturbation terms of second order. For an electron moving in an external electromagnetic field, such an equation was established in~\S\,33. Now we turn to deriving an analogous equation describing a system of two interacting particles.

We shall proceed from the relativistic expression for the amplitude of scattering of two particles. In the non-relativistic approximation it goes over to the ordinary Born scattering amplitude, proportional to the Fourier component of the potential of electrostatic interaction of two charges. Computing the amplitude to accuracy of terms of second order in~$1/c^2$, we shall be able to establish the form of the potential corresponding to it.

Let us suppose first that the two particles are different, with masses $m_1$ and $m_2$ (say, electron and muon). Then the scattering

\SourcePage{377}{382}
is depicted by a single diagram
% [Feynman diagram: Single photon exchange between two different fermion lines with momenta $p_1$, $p_1'$ and $p_2$, $p_2'$, photon momentum $q$]

To it there corresponds the amplitude
\begin{equation}
M_{fi} = e^2(\bar{u}_1'\gamma^\mu u_1)D_{\mu\nu}(q)(\bar{u}_2'\gamma^\nu u_2),\qquad q = p_1' - p_1 = p_2 - p_2'
\tag{83.1}
\end{equation}
(here it is assumed that the charges of the particles are of the same sign; in the opposite case $e^2$ is replaced by~$-e^2$).

Further calculations are substantially simplified if we choose the photon propagator $D_{\mu\nu}$ not in the ordinary but in the Coulomb gauge~(76.12),(76.13)\,\footnote{In this paragraph we write out the factors $c$ in all intermediate formulas, and in the final formulas also~$\hbar$.}:
\begin{equation}
D_{00} = -\frac{4\pi}{\mathbf{q}^2},\quad D_{0i} = 0,\quad D_{ik} = \frac{4\pi}{\mathbf{q}^2 - \omega^2/c^2}\left(\delta_{ik} - \frac{q_i q_k}{\mathbf{q}^2}\right).
\tag{83.2}
\end{equation}
Then the amplitude of scattering
\begin{equation}
M_{fi} = e^2\left\{(\bar{u}_1'\gamma^0 u_1)(\bar{u}_2'\gamma^0 u_2)D_{00} + (\bar{u}_1'\gamma^k u_1)(\bar{u}_2'\gamma^k u_2)D_{ik}\right\}.
\tag{83.3}
\end{equation}

To the neglect of all terms containing the factor $1/c$, the second member in the curly brackets drops out entirely, and the first gives
\begin{equation}
M_{fi} = -2m_1\cdot 2m_2\binom{w_1^{(0)'\ast}}{w_1^{(0)}}\binom{w_2^{(0)'\ast}}{w_2^{(0)}} U(\mathbf{q}),
\tag{83.4}
\end{equation}
where
\begin{equation}
U(\mathbf{q}) = \frac{4\pi e^2}{\mathbf{q}^2},
\tag{83.5}
\end{equation}
and by $w_1^{(0)}$, $w_2^{(0)}$, \ldots\ are denoted the 3-spinors of the non-relativistic plane waves introduced in~\S\,23 (two-component amplitudes). The function $U(\mathbf{q})$ represents the Fourier component of the potential energy of the Coulomb interaction: $U(r) = e^2/r$.

In the next (by $1/c$) approximation the ``Schr\"odinger'' wave function of the free particle $\varphi_\text{Schr}$ (normalised by the integral $\int|\varphi_\text{Schr}|^2 d^3x$) satisfies the equation
\begin{equation}
\hat{H}^{(0)}\varphi_\text{Schr} = (\varepsilon - mc^2)\varphi_\text{Schr},\qquad \hat{H}^{(0)} = \frac{\hat{\mathbf{p}}^2}{2m}-\frac{\hat{\mathbf{p}}^4}{8m^3c^2},\quad \hat{\mathbf{p}} = -i\hbar\boldsymbol{\nabla},
\tag{83.6}
\end{equation}
in which the next-order term of the relativistic expansion for the kinetic energy has been taken into account. The amplitude of the (spinor) part of such a plane wave we denote by $w$ (when $1/c \to 0$ it goes over to~$w^{(0)}$). It is precisely through these amplitudes that the desired scattering amplitude should be expressed so that from its form one could determine the ``Schr\"odinger'' interaction potential in the approximation being considered.

In accordance with formula~(33.11) the bispinor amplitude of the free particle $u$ is expressed through the ``Schr\"odinger'' amplitude $w$---with the accuracy required here---in the form
\begin{equation}
u = \sqrt{2m}\begin{pmatrix}\left(1 - \dfrac{\mathbf{p}^2}{8m^2c^2}\right)w\\[6pt] \dfrac{\boldsymbol{\sigma}\mathbf{p}}{2mc}w\end{pmatrix}.
\tag{83.7}
\end{equation}

\SourcePage{378}{383}
With the help of this formula we find
\[
\bar{u}_1'\gamma^0 u_1 = u_1'^*\alpha_1 u_1 = 2m_1\left(1-\frac{\mathbf{p}_1'^2+\mathbf{p}_1^2}{8m_1^2c^2}\right)w_1'^*w_1 +
\]
\[
+\frac{1}{2m_1c^2}w_1'^*(\boldsymbol{\sigma}\mathbf{p}_1')(\boldsymbol{\sigma}\mathbf{p}_1)w_1 = 2m_1 w_1'^*\left\{1-\frac{\mathbf{q}^2}{8m_1^2c^2}+\frac{i\boldsymbol{\sigma}[\mathbf{q}\mathbf{p}_1]}{4m_1^2c^2}\right\}w_1,
\]
\[
\bar{u}_1'\gamma u_1 = u_1'^*\boldsymbol{\alpha}_1 u_1 = \frac{1}{c}w_1'^*\{\boldsymbol{\sigma}(\boldsymbol{\sigma}\mathbf{p}_1) + (\boldsymbol{\sigma}\mathbf{p}_1')\boldsymbol{\sigma}\}w_1 =
\]
\[
= \frac{1}{c}w_1'^*\{i[\boldsymbol{\sigma}\mathbf{q}] + 2\mathbf{p}_1 + \mathbf{q}\}w_1
\]
(where $\mathbf{q} = \mathbf{p}_1' - \mathbf{p}_1 = \mathbf{p}_2 - \mathbf{p}_2'$). Analogous expressions for $(\bar{u}_2'\gamma^0 u_2)$ and $(\bar{u}_2'\gamma u_2)$ differ by the replacement of indices 1 by 2 and correspondingly the replacement of $q$ by~$-q$.

Substituting these expressions in~(83.3). Since the product $(\bar{u}_1'\gamma u_1)(\bar{u}_2'\gamma u_2)$ already contains the factor $1/c^2$, in it one can neglect the term $\omega^2/c^2$ in the denominator of $D_{ik}$. As a result we obtain the scattering amplitude in the form
\begin{equation}
M_{fi} = -2m_1\cdot 2m_2(w_1'^*w_2'^*U(\mathbf{p}_1, \mathbf{p}_2, \mathbf{q})\,w_1 w_2),
\tag{83.8}
\end{equation}
where
\begin{equation}
\begin{split}
U(\mathbf{p}_1, \mathbf{p}_2, \mathbf{q}) &= 4\pi e^2\bigg\{\frac{1}{\mathbf{q}^2} - \frac{1}{8m_1^2c^2} - \frac{1}{8m_2^2c^2} + \frac{(\mathbf{q}\mathbf{p}_1)(\mathbf{q}\mathbf{p}_2)}{m_1 m_2 c^2\mathbf{q}^4} - \frac{\mathbf{p}_1\mathbf{p}_2}{m_1 m_2 c^2\mathbf{q}^2}+{}\\
&\quad +\frac{i\boldsymbol{\sigma}_1[\mathbf{q}\mathbf{p}_1]}{4m_1^2c^2\mathbf{q}^2} - \frac{i\boldsymbol{\sigma}_1[\mathbf{q}\mathbf{p}_2]}{2m_1 m_2 c^2\mathbf{q}^2} - \frac{i\boldsymbol{\sigma}_2[\mathbf{q}\mathbf{p}_2]}{4m_2^2c^2\mathbf{q}^2} + \frac{i\boldsymbol{\sigma}_2[\mathbf{q}\mathbf{p}_1]}{2m_1 m_2 c^2\mathbf{q}^2}+{}\\
&\quad +\frac{(\boldsymbol{\sigma}_1\mathbf{q})(\boldsymbol{\sigma}_2\mathbf{q})}{4m_1 m_2 c^2\mathbf{q}^2} - \frac{\boldsymbol{\sigma}_1\boldsymbol{\sigma}_2}{4m_1 m_2 c^2}\bigg\}
\end{split}
\tag{83.9}
\end{equation}
(indices 1, 2 of the Pauli matrices indicate on which spinor indices they act: $\boldsymbol{\sigma}_1$ acts on $w_1$, and $\boldsymbol{\sigma}_2$ on~$w_2$).

\SourcePage{379}{384}
The function $U(\mathbf{p}_1, \mathbf{p}_2, \mathbf{q})$ is the operator of particle interaction in the momentum representation. It is related to the operator $\hat{U}(\hat{\mathbf{p}}_1, \hat{\mathbf{p}}_2, \mathbf{r})$ in the coordinate representation by the formula
\begin{equation}
\int e^{-i(\mathbf{p}_1'\mathbf{r}_1+\mathbf{p}_2'\mathbf{r}_2)}\hat{U}(\hat{\mathbf{p}}_1,\hat{\mathbf{p}}_2,\mathbf{r})\,e^{i(\mathbf{p}_1\mathbf{r}_1+\mathbf{p}_2\mathbf{r}_2)}\,d^3x_1\,d^3x_2 =
\end{equation}
\[
= (2\pi)^3\delta(\mathbf{p}_1+\mathbf{p}_2-\mathbf{p}_1'-\mathbf{p}_2')U(\mathbf{p}_1,\mathbf{p}_2,\mathbf{q}).
\tag{83.10}
\]
If the operator $\hat{U}$ represents simply a function $U(\mathbf{r})$ ($\mathbf{r} = \mathbf{r}_1 - \mathbf{r}_2$) that does not depend on $\mathbf{p}_1$ and $\mathbf{p}_2$, then formula~(83.10) reduces to the ordinary definition of the Fourier components:
\[
\int e^{-i\mathbf{qr}}U(\mathbf{r})\,d^3x = U(\mathbf{q}).
\]
From this it is clear that to find $\hat{U}(\hat{\mathbf{p}}_1, \hat{\mathbf{p}}_2, \mathbf{r})$ one must compute the integral
\[
\int e^{i\mathbf{qr}}U(\mathbf{p}_1, \mathbf{p}_2, \mathbf{q})\,\frac{d^3q}{(2\pi)^3}
\]
and then replace $\mathbf{p}_1$, $\mathbf{p}_2$ by operators $\hat{\mathbf{p}}_1 = -i\boldsymbol{\nabla}_1$, $\hat{\mathbf{p}}_2 = -i\boldsymbol{\nabla}_2$, placing them to the right of all other factors.

The necessary integrals are computed by differentiating the formula
\begin{equation}
\int e^{i\mathbf{qr}}\frac{4\pi}{\mathbf{q}^2}\,\frac{d^3q}{(2\pi)^3} = \frac{1}{r}.
\tag{83.11}
\end{equation}
Thus, taking the gradient we find
\begin{equation}
\int e^{i\mathbf{qr}}\cdot\frac{4\pi\mathbf{q}}{\mathbf{q}^2}\,\frac{d^3q}{(2\pi)^3} = -i\boldsymbol{\nabla}\frac{1}{r} = \frac{i\mathbf{r}}{r^3}.
\tag{83.12}
\end{equation}
Further ($\mathbf{a}$, $\mathbf{b}$ are constant vectors)
\begin{equation}
\int\frac{4\pi(\mathbf{aq})(\mathbf{bq})}{\mathbf{q}^4}\,e^{i\mathbf{qr}}\,\frac{d^3q}{(2\pi)^3} = \frac{i}{2}\left(\mathbf{a}\frac{\partial}{\partial\mathbf{r}}\right)\int e^{i\mathbf{qr}}\left(\mathbf{b}\frac{\partial}{\partial\mathbf{q}}\right)\frac{1}{\mathbf{q}^2}\,\frac{d^3q}{(2\pi)^3};
\tag{83.13}
\end{equation}
the resulting integral after integration by parts reduces to~(83.12) and gives
\[
\int\frac{4\pi(\mathbf{aq})(\mathbf{bq})}{\mathbf{q}^4}\,e^{i\mathbf{qr}}\,\frac{d^3q}{(2\pi)^3} = \frac{1}{2}(\mathbf{a}\boldsymbol{\nabla})\frac{\mathbf{br}}{r} = \frac{1}{2r}\left[\mathbf{ab}-\frac{(\mathbf{ar})(\mathbf{br})}{r^2}\right].
\]

Finally,
\[
\int\frac{4\pi(\mathbf{aq})(\mathbf{bq})}{\mathbf{q}^2}\,e^{i\mathbf{qr}}\,\frac{d^3q}{(2\pi)^3} = -(\mathbf{a}\boldsymbol{\nabla})(\mathbf{b}\boldsymbol{\nabla})\frac{1}{r}.
\]

\SourcePage{380}{385}
In expanding the derivatives one must bear in mind that this expression contains the $\delta$-function $\delta(\mathbf{r})$. To separate it we note that after averaging over the directions of~$\mathbf{r}$:
\[
-(\mathbf{a}\boldsymbol{\nabla})(\mathbf{b}\boldsymbol{\nabla})\frac{1}{r} = -\frac{1}{3}(\mathbf{ab})\Delta\frac{1}{r} = \frac{4\pi}{3}(\mathbf{ab})\delta(\mathbf{r}).
\]
Expanding now the derivatives in the ordinary way, we find
\begin{equation}
\int\frac{4\pi(\mathbf{aq})(\mathbf{bq})}{\mathbf{q}^2}\,e^{i\mathbf{qr}}\,\frac{d^3q}{(2\pi)^3} = \frac{1}{r^3}\left\{\mathbf{ab}-3\frac{(\mathbf{ar})(\mathbf{br})}{r^2}\right\} + \frac{4\pi}{3}(\mathbf{ab})\delta(\mathbf{r})
\tag{83.14}
\end{equation}
(on averaging over the directions of $\mathbf{r}$ the first term on the right vanishes and there remains only the term with the $\delta$-function).

With the help of these formulas we obtain the following final expression for the operator of particle interaction:
\begin{equation}
\begin{split}
\hat{U}(\hat{\mathbf{p}}_1, \hat{\mathbf{p}}_2, \mathbf{r}) &= \frac{e^2}{r} - \frac{\pi e^2\hbar^2}{2c^2}\left(\frac{1}{m_1^2}+\frac{1}{m_2^2}\right)\delta(\mathbf{r}) -{}\\
&- \frac{e^2}{2m_1 m_2 c^2 r}\left[\hat{\mathbf{p}}_1\hat{\mathbf{p}}_2+\frac{\mathbf{r}(\mathbf{r}\hat{\mathbf{p}}_1)\hat{\mathbf{p}}_2}{r^2}\right] - \frac{e^2\hbar}{4m_1^2c^2r^3}[\mathbf{r}\hat{\mathbf{p}}_1]\boldsymbol{\sigma}_1+\frac{e^2\hbar}{4m_2^2c^2r^3}[\mathbf{r}\hat{\mathbf{p}}_2]\boldsymbol{\sigma}_2-{}\\
&-\frac{e^2}{2m_1 m_2 c^2 r}\left\{-(\boldsymbol{\sigma}_1+2\boldsymbol{\sigma}_2)[\mathbf{r}\hat{\mathbf{p}}_1]+(\boldsymbol{\sigma}_2+2\boldsymbol{\sigma}_1)[\mathbf{r}\hat{\mathbf{p}}_2]\right\}+{}\\
&+ \frac{e^2\hbar^2}{4m_1 m_2 c^2}\left\{\frac{\boldsymbol{\sigma}_1\boldsymbol{\sigma}_2}{r^3}-3\frac{(\boldsymbol{\sigma}_1\mathbf{r})(\boldsymbol{\sigma}_2\mathbf{r})}{r^5} - \frac{8\pi}{3}\boldsymbol{\sigma}_1\boldsymbol{\sigma}_2\delta(\mathbf{r})\right\}.
\end{split}
\tag{83.15}
\end{equation}

The full Hamiltonian of the system of two particles in this approximation
\begin{equation}
\hat{H} = \hat{H}_1^{(0)} + \hat{H}_2^{(0)} + \hat{U},
\tag{83.16}
\end{equation}
where $\hat{H}^{(0)}$ are the Hamiltonians of free particles~(83.6).

\textbf{Two electrons.} If both particles are identical (two electrons), then in the scattering amplitude there appears a second term, depicted by the ``exchange'' diagram
% [Feynman diagram (exchange): Two crossed electron lines with momenta $p_1 \to p_2$, $p_2 \to p_1$, connected by a photon]

Computing its contribution to the interaction operator, however, is not necessary. The point is that the description of a system of identical particles by the Schr\"odinger equation can be accomplished with the help of the same interaction operator as for non-identical particles, if one takes care of the proper symmetrisation of the

\SourcePage{381}{386}
wave functions. In particular, in the consideration of the scattering of identical particles such symmetrisation automatically takes into account the contributions to the amplitude corresponding to both Feynman diagrams.

Thus, the Hamiltonian of a system of two electrons is obtained from formulas~(83.15),\,(83.16), if one simply puts $m_1 = m_2$\,\footnote{The wave equation with Hamiltonian~(83.17) was first established by \textit{Breit} (\textit{G.\,Breit}, 1929), and its consistent quantum-mechanical derivation was given by \textit{L.\,D.\,Landau}~(1932).}:
\begin{equation}
\hat{H} = \frac{1}{2m}(\hat{\mathbf{p}}_1^2+\hat{\mathbf{p}}_2^2) - \frac{1}{8m^3c^2}(\hat{\mathbf{p}}_1^4+\hat{\mathbf{p}}_2^4) + \hat{U}(\hat{\mathbf{p}}_1, \hat{\mathbf{p}}_2, \mathbf{r}),
\end{equation}
\begin{equation}
\begin{split}
\hat{U}(\hat{\mathbf{p}}_1, \hat{\mathbf{p}}_2, \mathbf{r}) &= \frac{e^2}{r} - \pi\left(\frac{e\hbar}{mc}\right)^2\delta(\mathbf{r}) - \frac{e^2}{2m^2c^2r}\left(\hat{\mathbf{p}}_1\hat{\mathbf{p}}_2+\frac{\mathbf{r}(\mathbf{r}\hat{\mathbf{p}}_1)\hat{\mathbf{p}}_2}{r^2}\right)+{}\\
&+\frac{e^2\hbar}{4m^2c^2r^3}\{-(\boldsymbol{\sigma}_1+2\boldsymbol{\sigma}_2)[\mathbf{r}\hat{\mathbf{p}}_1]+(\boldsymbol{\sigma}_2+2\boldsymbol{\sigma}_1)[\mathbf{r}\hat{\mathbf{p}}_2]\}+{}\\
&+\frac{1}{4}\left(\frac{e\hbar}{mc}\right)^2\frac{1}{r^3}\left\{\frac{\boldsymbol{\sigma}_1\boldsymbol{\sigma}_2}{r^3}-3\frac{(\boldsymbol{\sigma}_1\mathbf{r})(\boldsymbol{\sigma}_2\mathbf{r})}{r^5}-\frac{8\pi}{3}\boldsymbol{\sigma}_1\boldsymbol{\sigma}_2\delta(\mathbf{r})\right\}.
\end{split}
\tag{83.17}
\end{equation}

Let us note that the presence of terms with $\delta(\mathbf{r})$ does not, of course, signify the existence of a special strong interaction. The integral magnitude of all correction terms is the same, and in the spirit of the expansion performed they must all be considered small in comparison with the first term---the Coulomb interaction.

The various groups of terms in the interaction operator~(83.17) have different character. The terms of the first row in $\hat{U}$ have a purely orbital origin. In the second row stand the terms, linear in the spin operators of the particles; they correspond to the \textit{spin-orbit interaction}. Finally, the terms of the third row, quadratic in the spin operators, describe the \textit{spin-spin interaction}\,\footnote{This interaction was mentioned in~III, \S\,72 in connection with the fine structure of atomic levels, and the spin-spin interaction of electrons with the nucleus was treated in~III, \S\,121, in connection with the superhyperfine structure of levels. In particular, formula~(121.9) (see~III) corresponds to the $\delta$-function term in the spin-spin interaction operator.}.

\textbf{Electron and positron.} The system of electron and positron requires special consideration. The scattering amplitude in this case is composed of two terms:
\begin{equation}
\begin{split}
M_{fi} &= -e^2[\bar{u}(p_-')\gamma^\mu u(p_-)]D_{\mu\nu}(p_- - p_-')[\bar{u}(-p_+)\gamma^\nu u(-p_+')]+{}\\
&+ e^2[\bar{u}(-p_+)\gamma^\mu u(p_-)]D_{\mu\nu}(p_- + p_+)[\bar{u}(p_-')\gamma^\nu u(-p_+')]
\end{split}
\tag{83.18}
\end{equation}
(the first corresponds to the scattering diagram, and the second to the annihilation diagram). Since the wave function of the ``electron +

\SourcePage{382}{387}
positron'' system should not be antisymmetric, both terms give independent contributions to the operator of interaction.

The first term (whose structure coincides with the structure of the scattering amplitude~(83.1)) leads, naturally, to an operator differing from~(83.17) only by the overall sign.

Let us turn to the transformation of the second term.

We use here the photon propagator in the ordinary gauge:
\begin{equation}
D_{\mu\nu} = \frac{4\pi}{k^2}g_{\mu\nu} = \frac{4\pi}{\omega^2/c^2 - \mathbf{k}^2}g_{\mu\nu}.
\tag{83.19}
\end{equation}

In the present case $k = p_+ + p_-$, and since the particles are ``almost non-relativistic'', we have
\[
\frac{\omega^2}{c^2} \equiv \frac{(\varepsilon_+ + \varepsilon_-)^2}{c^2} \approx 4m^2c^2 \gg (\mathbf{p}_+ + \mathbf{p}_-)^2 \equiv \mathbf{k}^2.
\]
Therefore for the photon propagator it suffices to write
\[
D_{\mu\nu} \approx \frac{\pi}{m^2c^2}g_{\mu\nu}.
\]
This already contains the factor $1/c^2$. Therefore the amplitudes $u(p)$ can be taken in the zeroth approximation:
\[
u(p_-) = \sqrt{2m}\begin{pmatrix}w_-^{(0)}\\0\end{pmatrix},\qquad u(-p_+) = \sqrt{2m}\begin{pmatrix}0\\w^{(0)}\end{pmatrix},
\]
where $w_-^{(0)}$, $w^{(0)}$ are the 3-spinors figuring in~(23.12) (below we omit the indices~(0) on them). With these amplitudes
\[
\bar{u}(-p_+)\gamma^0 u(p_-) = u^*(-p_+) u(p_-) = 0,
\]
\[
\bar{u}(-p_+)\gamma u(p_-) = u^*(-p_+)\boldsymbol{\alpha}u(p_-) = 2m(w^*\boldsymbol{\sigma}w_-),
\]
where $w_-^{(0)}$, $w^{(0)}$ are the 3-spinors figuring in~(23.12). With these amplitudes
\begin{equation}
M_{fi}^\text{(ann)} = -e^2\frac{\pi}{m^2c^2}(2m)^2(w^*\boldsymbol{\sigma}w_-)(w'^{*}_-\boldsymbol{\sigma}w').
\tag{83.20}
\end{equation}
From this, however, one cannot yet draw direct conclusions about the form of the operator of interaction. First, the scattering amplitude should be brought to the form in which the electronic ($w_-$ and $w_-'$) and positronic ($w_+$ and $w_+'$) spinors are grouped together. This is achieved with the help of the formula
\begin{equation}
(w^*\boldsymbol{\sigma}w_-)(w'^{*}_-\boldsymbol{\sigma}w') = \frac{3}{2}(w'^*_-w_-)(w^*w') - \frac{1}{2}(w'^*_-\boldsymbol{\sigma}w_-)(w^*\boldsymbol{\sigma}w'),
\tag{83.22}
\end{equation}
which itself follows from~(28.16).

Finally, expressing $w$ and $w'$ through $w_+$ and $w_+'$, according to~(83.21), we find, as is easy to verify,
\begin{equation}
(w^*w') = (w_+'^*w_+),\qquad (w^*\boldsymbol{\sigma}w') = -(w_+'^*\boldsymbol{\sigma}w_+).
\tag{83.23}
\end{equation}

\SourcePage{383}{388}
Substituting~(83.23) in~(83.22) and then in~(83.20), we obtain the final expression for the annihilation part of the scattering amplitude
\begin{equation}
M_{fi}^\text{(ann)} = -4m^2\left\{w'^*_- w'^*_+\left[\frac{\pi e^2}{2m^2c^2}(3+\boldsymbol{\sigma}_+\boldsymbol{\sigma}_-)\right]w_- w_+\right\}
\end{equation}
(the matrices $\boldsymbol{\sigma}_-$ and $\boldsymbol{\sigma}_+$ act respectively on $w_-$ and~$w_+$). The expression in the square brackets represents the interaction operator in the momentum representation. The corresponding coordinate operator
\begin{equation}
\hat{U}^\text{(ann)}(\mathbf{r}) = \frac{\pi\hbar^2 e^2}{2m^2c^2}(3+\boldsymbol{\sigma}_+\boldsymbol{\sigma}_-)\delta(\mathbf{r}),\qquad \mathbf{r} = \mathbf{r}_- - \mathbf{r}_+
\tag{83.24}
\end{equation}
(\textit{Pirenne}, 1947; \textit{B.\,B.\,Berestetskii} and \textit{L.\,D.\,Landau}, 1949). The full interaction operator of the electron and positron is
\[
-\hat{U} + \hat{U}^\text{(ann)}
\]
with $\hat{U}$ from~(83.17).
